% !TeX root = ../informe.tex
% ===========================================================================
%  7. Resultados
%
%  Cierra el círculo con los objetivos (sección 2) y con los criterios de
%  aceptación del documento de alcance. Los \label{ca:N} son los que hacen
%  funcionar \crit{N} desde el resto del documento.
% ===========================================================================

\section{Resultados}
\label{sec:resultados}

\subsection{Funcionalidades implementadas}
\label{sec:funcionalidades}

\instruccion{Qué quedó operativo, con capturas de las pantallas principales.
Dos o tres capturas bien elegidas valen más que ocho: la consulta de
disponibilidad, la reserva y los reportes.}

A la fecha de este informe están construidas y demostrables tres de las cuatro
historias del alcance:

\begin{itemize}
  \item \textbf{Reservar y cancelar como usuario final.} Consulta de
        disponibilidad por cancha y fecha, creación de reserva sobre un bloque
        libre, listado propio y cancelación que libera el horario.
  \item \textbf{Gestión del catálogo y cancelación administrativa.} Alta,
        edición, activación e inactivación de canchas, bloqueos de
        mantenimiento, listado global de reservas y cancelación de cualquiera.
  \item \textbf{Los cuatro indicadores de §3.3.5}, sobre el rango que el
        administrador elija, más un panel acotado al día en curso.
\end{itemize}

Queda pendiente la gestión de usuarios —activar e inactivar cuentas desde la
pantalla de administración—, cuyo respaldo en \id{ms-usuarios} sí existe: el
rechazo de una cuenta inactiva al iniciar sesión está operativo y probado.

Toda la infraestructura está en pie: \id{docker compose up} levanta once
contenedores —edge, cuatro microfrontends, gateway, cuatro microservicios y
PostgreSQL con sus tres bases aisladas— con datos de prueba y sin ningún paso
manual.

\subsection{Cumplimiento de los criterios de aceptación}
\label{sec:cumplimiento}

\instruccion{Los seis criterios del documento de alcance, uno por fila, con
la evidencia que respalda cada «Sí». Es la tabla que un evaluador busca
primero. Si alguno no se cumple, decláralo: un criterio marcado como
cumplido que luego falla en la demo cuesta más que uno declarado pendiente.}

\begin{table}[H]
  \centering
  \caption{Cumplimiento de los criterios de aceptación.}
  \label{tab:cumplimiento}
  \begin{tabularx}{\textwidth}{C{1.4cm} Y C{2cm} L{3cm}}
    \toprule
    \textbf{\#} & \textbf{Criterio} & \textbf{Estado} & \textbf{Evidencia} \\
    \midrule
    \textsf{CA-1}\label{ca:1} & Consultar, crear y cancelar reservas de las tres canchas & Cumplido & Escenario V1 del \id{quickstart}; 24 pruebas de \id{BookingService}. \\
    \textsf{CA-2}\label{ca:2} & El administrador gestiona el catálogo y cancela cualquier reserva & Cumplido & Escenario V3; 13 pruebas de \id{CourtService} y la rama de administrador en \id{BookingServiceCancelacionTest}. \\
    \textsf{CA-3}\label{ca:3} & El sistema impide reservar un bloque ocupado & Cumplido & Escenario V2, diez repeticiones: siempre un 201 y un 409. \id{ConcurrenciaReservaIT} contra PostgreSQL real (\cref{sec:solapamiento}). \\
    \textsf{CA-4}\label{ca:4} & Los reportes muestran ocupación y reservas por período & Cumplido & Escenario V4: los cuatro indicadores contrastados contra un conteo manual del listado global del mismo rango. \\
    \textsf{CA-5}\label{ca:5} & Shell y al menos dos remotes integrados con Module Federation & Cumplido & Tres remotes —reservas, administración y reportes— cargados de forma perezosa por el shell (\cref{v:02}). \\
    \textsf{CA-6}\label{ca:6} & Al menos dos microservicios independientes sobre PostgreSQL & Cumplido & Tres con base propia y aislada por credenciales, más \id{ms-reportes} sin base (\cref{sec:estrategia-datos}). \\
    \bottomrule
  \end{tabularx}
\end{table}

\subsection{Verificación de las reglas de negocio}
\label{sec:verificacion-reglas}

\instruccion{Cómo se comprobó que cada regla funciona. No hace falta un
capítulo de pruebas: una tabla con el escenario, el resultado esperado y el
obtenido es suficiente y se lee en treinta segundos. La colección de
peticiones va al anexo.}

\begin{table}[H]
  \centering
  \caption{Comprobación de las reglas de negocio.}
  \label{tab:pruebas-reglas}
  \begin{tabular}{C{1.6cm} L{6.6cm} L{4.8cm}}
    \toprule
    \textbf{Regla} & \textbf{Escenario} & \textbf{Resultado} \\
    \midrule
    \rn{01} & Reservar cancha 1, mañana, bloque 18:00. & 201, bloque 18:00--19:00 confirmado. \\
    \rn{02} & Un segundo usuario pide el mismo bloque. & 409 con el motivo en pantalla. \\
    \rn{02} & Dos peticiones simultáneas, diez repeticiones. & Siempre un 201 y un 409; nunca dos. \\
    \rn{03} & \id{cliente1} cancela una reserva de \id{cliente2}. & 403. El administrador sobre la misma: 200. \\
    \rn{04} & Cancelar una reserva ya iniciada. & 409, también siendo administrador. \\
    \rn{05} & Cancelar una reserva futura y reconsultar. & El bloque vuelve a figurar libre. \\
    \rn{06} & Crear una cuarta reserva activa con el tope en 3. & 409 indicando el máximo. \\
    \rn{07} & \id{cliente1} intenta crear una cancha. & 403; el administrador, 201. \\
    \rn{08} & Listar las reservas propias. & Los tres estados, con \id{FINALIZADA} derivada al leer. \\
    \bottomrule
  \end{tabular}
\end{table}

\subsection{Limitaciones}
\label{sec:limitaciones}

\instruccion{Qué no se alcanzó y por qué. Declararlo suma: demuestra que se
conoce el propio sistema. Omitir una limitación que el evaluador descubre en
la demo resta el doble.}

\begin{itemize}
  \item \textbf{La gestión de usuarios no está conectada a la pantalla.} El
        backend distingue cuentas activas e inactivas y rechaza el acceso de
        estas últimas, pero \id{mf-administracion} todavía no consume
        \id{/api/usuarios}.
  \item \textbf{No hay paginación en ningún listado.} A escala de un club los
        conjuntos son pequeños, y añadirla sería complejidad sin criterio de la
        rúbrica que la respalde.
  \item \textbf{\id{ms-reportes} agrega en memoria}, no con consultas agregadas.
        Es el precio de no leer la base de otro servicio
        (\cref{sec:estrategia-datos}), y es un precio consciente.
  \item \textbf{Las pruebas de contexto de Spring necesitan una base real} y por
        eso no se ejecutan en un entorno sin PostgreSQL. Las pruebas de las
        reglas de negocio, que son las que el Principio~II exige, corren siempre:
        no levantan Spring ni base de datos.
\end{itemize}
